\documentclass[11pt]{article}
\usepackage[a4paper,margin=25mm]{geometry}
\usepackage{amsmath,amssymb,amsthm,hyperref}
\newtheorem{lemma}{Lemma}\newtheorem{theorem}{Theorem}
\newcommand{\E}{\mathbb E}
\title{Positive finite Jackson kernels at the endpoint-profile scale}
\author{Complete elementary analytic argument; independently reviewed}
\date{30 September 2026}
\begin{document}\maketitle
\noindent\textbf{Status.} The complete note passed independent mathematical reviews by the root, independent-directions, and size-bounds agents. This is not formal verification.

The kernels in this note are actual polynomials of degree $O(m)$ on
$[0,\ell]$. Their interior endpoint-profile width is $m^{-1/2}$;
there is no infinite spectral tail. The global polarization tail is a
separate Bernstein--Walsh estimate.

Fix an integer $r\ge2$. For an integer $\nu\ge2$, put
\[
 F_\nu(t)=\left(\frac{\sin(\nu t/2)}{\sin(t/2)}\right)^2,
 \quad Z_{\nu,r}=\frac1{2\pi}\int_{-\pi}^{\pi}F_\nu(t)^r\,dt,
 \quad J_{\nu,r}(t)=F_\nu(t)^r/Z_{\nu,r}.
\]
At removable singularities use the continuous values. This is a
nonnegative even trigonometric polynomial of degree $r(\nu-1)$.
Write $|t|_{2\pi}=\min_{j\in\mathbb Z}|t-2\pi j|$.

\begin{lemma}\label{trig}
With constants depending only on $r$,
\[
 Z_{\nu,r}\asymp\nu^{2r-1},\qquad
 0\le J_{\nu,r}(t)\le
       \frac{C_r\nu}{(1+\nu|t|_{2\pi})^{2r}},
\]
and $J_{\nu,r}(t)\ge c_r\nu$ if $|t|_{2\pi}\le c_r/\nu$.
For real $t$ and complex $\eta$,
\[
 |J_{\nu,r}(t+\eta)|\le
 \frac{C_r\nu\exp(C_r\nu|\eta|)}
      {(1+\nu|t|_{2\pi})^{2r}}.                         \tag{1}
\]
\end{lemma}
\begin{proof}
The quotient of sines is a phased sum of $\nu$ exponentials.
For $w=u+iv$ its absolute value is at most
\[
 C\min(\nu,|u|_{2\pi}^{-1})e^{\nu|v|/2}
 \le\frac{C\nu e^{\nu|v|/2}}{1+\nu|u|_{2\pi}}.
\]
The first bound follows both from the exponential sum and from the
lower bound on $|\sin(w/2)|$ in terms of $|u|_{2\pi}$.
For real $|t|\le c/\nu$ the quotient is at least $c\nu$.
Integrating the upper and lower bounds proves the normalization claim
and the real estimates. Finally
\[
 1+\nu|t|_{2\pi}\le
 (1+\nu|\operatorname{Re}\eta|)
 (1+\nu|t+\operatorname{Re}\eta|_{2\pi});
\]
the resulting fixed power of $1+\nu|\eta|$ is absorbed into the
exponential. This proves (1).
\end{proof}

Assume throughout that $c_1m\le\ell\le c_2m$ and
$c_3m\le\nu\le c_4m$, with positive constants fixed before $m$.
The upper constant $c_4$ can be chosen arbitrarily small when needed
for the global polarization estimate. For $x\in[0,\ell]$, define
\[
 \theta(x)=\arccos(1-2x/\ell),\qquad b=\sqrt\ell/\nu.
\]
For a center $a$ in a fixed compact interval $[\delta,L]\subset(0,\infty)$,
put $\theta_a=\theta(a)$ and
\[
 Q_m(x,a)=\frac{J_{\nu,r}(\theta(x)-\theta_a)
                   +J_{\nu,r}(\theta(x)+\theta_a)}
                  {\ell\sin\theta_a}.                  \tag{2}
\]

\begin{theorem}[Positive polynomial localization]\label{kernel}
For all sufficiently large $m$, $Q_m(\cdot,a)$ is a nonnegative
ordinary polynomial on $[0,\ell]$ of degree at most $r(\nu-1)$,
and $\|Q_m(\cdot,a)\|_{\infty,[0,\ell]}\le C/b$.
On any fixed positive compact interval of $x$ containing the centers,
\[
 0\le Q_m(x,a)\le C E_b(x,a),\qquad
 E_b(x,a)=\frac1b(1+|x-a|/b)^{-2r}.                       \tag{3}
\]
Also $Q_m(x,a)\ge c/b$ whenever $|x-a|\le cb$.
There are fixed $0<x_-<x_+<\infty$ such that
\[
 Q_m(x,a)\le C b^{2r-1}(1+x)^{-r}
 \quad\text{for }x\in[0,\ell]\setminus[x_-,x_+].         \tag{4}
\]
If $\beta_\ell$ has density
$\frac{\ell+1}{\ell}(1-x/\ell)^\ell$ on $[0,\ell]$, then
\[
 c\le\int Q_m(x,a)\,d\beta_\ell(x)\le C.
\]
All constants are uniform in $a\in[\delta,L]$; they may depend on
$r$ and the fixed ratio and interval constants.
\end{theorem}
\begin{proof}
Writing the even trigonometric polynomial as a cosine series shows
that the numerator of (2) is a linear combination of
$\cos(j\theta(x))=T_j(1-2x/\ell)$, since
$\cos(j(\theta-\theta_a))+\cos(j(\theta+\theta_a))
=2\cos(j\theta)\cos(j\theta_a)$. This proves polynomiality and degree;
nonnegativity is immediate. Also
$\ell\sin\theta_a=2\sqrt{a(\ell-a)}\asymp\sqrt\ell$.
The supremum estimate follows from Lemma~\ref{trig}.

On fixed positive compact intervals,
$|\theta(x)-\theta(a)|\asymp|x-a|/\sqrt\ell$.
The reflected angle $\theta(x)+\theta_a$ is bounded below by
$c/\sqrt\ell$ and above by $C/\sqrt\ell$.
Its contribution is $O(b^{2r-1})$, which is absorbed into (3) on a
fixed interval. The first angle gives (3) and the near-diagonal lower
bound. If $x$ is below a sufficiently small fixed $x_-$, both angles
stay a fixed multiple of $\ell^{-1/2}$ from zero. If $x$ is above a
sufficiently large fixed $x_+$, both periodic distances are bounded
below by $c\sqrt{x/\ell}$; at the upper end $x\approx\ell$ they
are uniformly separated from $0$ and $2\pi$.
Lemma~\ref{trig} then proves (4).
The reference density is bounded above and below on fixed compact
intervals and is at most $2e^{-x}$. Integrate (3)--(4), using the
lower bound on a $cb$ neighborhood of $a$, to finish.
\end{proof}

\begin{lemma}[Localized complex derivatives]\label{derivatives}
Fix a positive compact interval of $x$ and centers $a$.
For some fixed $c>0$, all complex $|\zeta|\le c$ satisfy
\[
 |Q_m(x+\zeta,a)|\le C E_b(x,a)e^{C|\zeta|/b}.
\]
Consequently for every integer $k\ge1$,
\[
 \frac{|\partial_x^k Q_m(x,a)|}{k!}
 \le C E_b(x,a)\left[\left(\frac C{kb}\right)^k+C^k\right].
                                                                    \tag{5}
\]
\end{lemma}
\begin{proof}
The branch $\theta(z)$ is analytic in a fixed complex neighborhood
of the positive compact interval, and
$\theta'(z)=1/\sqrt{z(\ell-z)}$. Hence
$|\theta(x+\zeta)-\theta(x)|\le C|\zeta|/\sqrt\ell$ there.
Apply (1) to both angles in (2), and use the same envelope comparison
as in (3). This proves the complex bound.
Cauchy's formula on the radius
$R=\min(c,kb/C_0)$, with $C_0$ sufficiently large, proves (5).
In the capped-radius case $k\ge c'/b$, so the exponential
$e^{Cc/b}$ is at most $e^{C'k}$ and is absorbed into $C^k$.
\end{proof}

\begin{lemma}[Finite centered differential sums]\label{polar}
Suppose the row values $z_j$ and center $a$ lie in a fixed positive
compact interval, let $v$ be the mean of $N\asymp m$ of those values,
and let $H$ be their total range, including a marked value $z_i$.
Set $a_k=e_k((z_j-v)_{j\in R})/\binom Nk$.
For $H$ sufficiently small and every $D\le N$,
\begin{align*}
 \left|\sum_{k=2}^D\frac{Q_m^{(k)}(v,a)}{k!}a_k\right|
 &\le C E_b(v,a)\frac{H^2}{mb^2},\\
 \left|\sum_{k=2}^D
 \left[\frac{(v-z_i)Q_m^{(k)}(v,a)}{k!}
       +\frac{Q_m^{(k-1)}(v,a)}{(k-1)!}\right]a_k\right|
 &\le C E_b(v,a)
             \left(\frac{H^2}{mb}+\frac{H^3}{mb^2}\right).
\end{align*}
The estimates remain valid after success-weighted subset averaging,
with $v$ replaced by its unweighted population mean $x$ in the
envelope. Replacing the scalar quantities $Q_m(v,a)$ and
$(z_i-v)Q_m(v,a)$ by $Q_m(x,a)$ and $(z_i-x)Q_m(x,a)$ costs at most
\[
 C E_b(x,a)\frac{H^2}{mb^2},\qquad
 C E_b(x,a)\frac{H^2}{m}(1+b^{-1}+Hb^{-2}),               \tag{6}
\]
respectively.
\end{lemma}
\begin{proof}
Use $|a_k|\le(CH\sqrt{k/N})^k$. The first term in (5) gives a series
bounded by
$\sum_{k\ge2}(CH/(b\sqrt{Nk}))^k=O(H^2/(Nb^2))$.
The $C^k$ term gives $O(H^2/N)$, since
$(k/N)^{k/2}\le k/N$ for $2\le k\le N$ and $H$ is small.
For derivative order $k-1$ the first series is at most
\[
 \frac{CH}{\sqrt N}
    \sum_{k\ge2}\sqrt{k}
             \left(\frac{CH}{b\sqrt N}\right)^{k-1}
 =O(H^2/(Nb));
\]
this deliberately weakens the extra powers of $k$ in the denominator.
The second series contributes $O(H^2/N)$.
Multiplying the order-$k$ part by $H$ proves the two fixed-subset bounds.

For the weighted subset law, the reviewed all-tilt lemma gives
$Y=v-x$ with bias $O(H^2/m)$ and centered subgaussian variance proxy
$CH^2/m$. Since $b\asymp m^{-1/2}$ and $H\to0$, it follows that
\[
 \E e^{C|Y|/b}\le C',\qquad
 \E[Y^2e^{C|Y|/b}]\le C'H^2/m.
\]
For instance use $e^{C|Y|/b}\le C_\epsilon e^{\epsilon Y^2/b^2}$
and the squared-exponential estimate in that lemma.
Also $E_b(v,a)\le E_b(x,a)(1+|Y|/b)^{2r}$, so averaging preserves
the envelope. Finally the localized complex bound and Cauchy's
formula at orders one and two give the derivative envelopes needed
for Taylor's theorem at $x$. The bias and the weighted second
moment just displayed give (6), exactly as for the reviewed Poisson
kernel comparison. All constants may depend on the fixed $r$.
\end{proof}


\begin{lemma}[Low-order derivatives away from the center]\label{far}
Fix the positive center interval and constants above, and let
$d=r(\nu-1)$. For every fixed $B$, all sufficiently large $m$,
$1\le k\le B\log m$, and
$v\in(0,\ell)\setminus[x_-,x_+]$, the constants $x_-,x_+$ in
Theorem~\ref{kernel} may be chosen so that
\[
 \frac{|Q_m^{(k)}(v,a)|}{k!}
 \le C_r b^{2r-1}(1+v)^{-r}
          \left(\frac{C_r d}{k\sqrt{v(\ell-v)}}\right)^k. \tag{7}
\]
\end{lemma}
\begin{proof}
Put $R=k\sqrt{v(\ell-v)}/d$. On the complex disk $|z-v|\le R$
there is a choice of the inverse-cosine value satisfying
\[
 |\theta(z)-\theta(v)|\le Ck/d.                         \tag{8}
\]
This is a bound on values; no analytic branch across an endpoint is
required, since the symmetric expression defining $Q_m$ is already
an ordinary polynomial. Here are details of (8). If
$R\le\min(v,\ell-v)/8$, integrate
$\theta'(z)=1/\sqrt{z(\ell-z)}$ to get
$C R/\sqrt{v(\ell-v)}$. Otherwise use the uniform square-root
continuity bound for the inverse cosine, giving $C\sqrt{R/\ell}$.
That bound follows by integrating its inverse-square-root derivative
near either endpoint; the integral of $|z-z_0|^{-1/2}$ along a segment
of length $h$ is at most $C\sqrt h$. Away from the endpoints the
ordinary derivative is bounded. In the second case, if
$a_0=\min(v,\ell-v)$, the inequality $R>a_0/8$ implies
$a_0/\ell\le C(k/d)^2$, whence $R/\ell\le C(k/d)^2$.
This proves (8). For $k\le B\log m$, the disks have relative radius
$R/\ell=o(1)$, so all the stated elementary inverse-cosine bounds are uniform.

Apply the complex estimate (1) to both angles defining $Q_m$.
By (8) its exponential cost is at most $e^{C_r k}$.
The real-angle envelopes away from the center were already bounded
in (4) by $C_rb^{2r-1}(1+v)^{-r}$.
Cauchy's formula on radius $R$ now gives (7).
\end{proof}

\begin{theorem}[Global off-annulus differential tails]\label{offannulus}
Use the exact marked-tensor notation and the two global inputs
\[
 \max_jz_j\le C(s+1),\qquad
 \E_w e^{-cs}\le C_c/m\quad(c>0\text{ fixed}),
\]
from \path{global_laguerre_spectral_truncation.tex}.
Let $a_k$ be the centered complementary coefficients and
$J=\lceil B\log m\rceil$ for any fixed $B$.
Fix a desired exponent $A>0$. If the fixed integer $r$ is sufficiently
large in terms of $A$, and then $d/m$ is a sufficiently small fixed
positive constant, there is a fixed annulus $[s_-,s_+]$ such that
\begin{align*}
 &(\ell+1)\E_w\E_B A_B\,\boldsymbol1_{s\notin[s_-,s_+]}\\
 &\quad\times\left[(1+|z_i-v|)|Q_m(v,a)|+
 \sum_{k=2}^{J}|a_k|\left\{
 (1+|z_i-v|)\frac{|Q_m^{(k)}(v,a)|}{k!}
       +\frac{|Q_m^{(k-1)}(v,a)|}{(k-1)!}\right\}\right]
 \le m^{-A}.
\end{align*}
This covers both unmarked and marked low-order differential sums.
\end{theorem}
\begin{proof}
Choose $s_-$ small enough that $v\le\ell s/N\le x_-$ on its rows,
and $s_+$ large enough that $x=(\ell/M)s$ satisfies $x/2\ge x_+$
on its rows. On large rows first consider subsets with $v\ge x/2$.
When $v\le\ell/2$, combine (7) with
\[
 |a_k|\le[Ck(s+1)v/N]^{k/2}
\]
to obtain the same integrated estimate as in the reviewed global
Bernstein--Walsh tail note, now with the additional factor
$b^{2r-1}$. More explicitly, after reserving the success exponential
$\E_BA_B\le e^{-\ell s/M}$ and using the Laplace row moments, the
$k$th normalized term is at most
\[
 C_rb^{2r-1}(k+1)^3(C_r d/m)^k
\]
for derivative order $k$, and at most
$C_rb^{2r-1}(k+1)^3m^{-1/2}(C_r d/m)^{k-1}$ for order $k-1$.
The scalar term has the same $O_r(b^{2r-1})$ bound.
The argument also applies to every small row, for which the far
condition is deterministic. The cases $v=0$ contribute only the
scalar term because all $a_k$, $k\ge1$, vanish.
Choose $d/m$ small enough to sum these estimates. Their total is
$O_r(b^{2r-1})=O_r(m^{1/2-r})$.

On subsets with $v>\ell/2$, use the global estimate from
\path{../../bernstein_walsh_polarization_tail.tex}:
$\sigma^2\le Nv(\ell-v)$ bounds all derivative terms by a polynomial
factor times $\|Q_m\|_\infty\exp(Cd^2/N)$.
Such rows have $s>N/2$, and the averaged success product is at most
$e^{-cm}$. For sufficiently small $d/m$ this contribution is
exponentially small. This applies regardless of whether $v\ge x/2$.

It remains to handle large rows with $v<x/2$ and $v\le\ell/2$.
As proved by the uniform-subset Maclaurin--Hoeffding argument in the
global truncation note, their unweighted subset probability is at
most $e^{-cm}$. A crude derivative bound suffices: the global
Bernstein--Walsh square-root estimate, followed by Cauchy's formula
at radius $\ell/d^2$, gives
\[
 |Q_m^{(k)}(v,a)|/k!\le
 C\|Q_m\|_\infty(Cd^2/\ell)^k.
\]
All row values are $O(m)$, and $k\le J=O(\log m)$, so the entire
finite sum, including the coefficients and marked factors, is at
most $\exp(C_B(\log m)^2)$. Since $A_B\le1$, its exceptional-subset
contribution is negligible. Finally choose $r>A+2$ (or increase it
to absorb any fixed polynomial prefactors required by an application)
to make $O_r(m^{1/2-r})\le\tfrac12m^{-A}$ for large $m$.
The exponentially small terms give the other half.
\end{proof}

\paragraph{Scope.}
These lemmas provide positive polynomial tests, reference mass,
compact differential comparison, and global off-annulus control at
width $m^{-1/2}$. The global high-differential-order estimate is proved
in the separate Bernstein--Walsh note cited above. The final tensor
application still requires the localized mass and weighted maximum
arguments with polynomial rather than Gaussian tails. No final profile
theorem is claimed in this note.
\end{document}
